\documentclass[10pt]{article}
\usepackage[T1,T2A]{fontenc}
\usepackage[utf8]{inputenc}
\usepackage[english,ukrainian]{babel}
\usepackage{geometry}
\usepackage{longtable}
\usepackage{array}
\usepackage{multirow}
\usepackage{hyperref}
\usepackage{mathptmx}
\renewcommand{\familydefault}{\rmdefault}

\geometry{a4paper,left=2cm,right=2cm,top=2cm,bottom=2cm}

\begin{document}

\begin{flushright}
ЗАТВЕРДЖЕНО\\
Наказом ТОВ "Криптографічні Телесистеми"\\
від \_\_ \_\_\_\_\_\_\_\_\_ 2026 р. № \_\_\_\\
(в редакції наказу ТОВ "Криптографічні Телесистеми"\\
від \_\_ \_\_\_\_\_\_\_\_\_ 2026 р. № \_\_)
\end{flushright}

\vspace{2cm}

\begin{center}
\textbf{\Large Flutter based X.422.1 Implementation}\\
\textbf{\Large «Chat X.422.1 (Flutter)»}\\
\vspace{1cm}
\textbf{\Large ЦІЛЬОВИЙ ПРОФІЛЬ БЕЗПЕКИ}\\
\vspace{0.5cm}
\large UA.47850061.СЗІ.ДП-50
\end{center}

\newpage

\footnotesize
\begin{longtable}{|>{\centering\arraybackslash}p{0.5cm}|>{\raggedright\arraybackslash}p{4cm}|>{\raggedright\arraybackslash}p{5cm}|>{\raggedright\arraybackslash}p{2cm}|>{\raggedright\arraybackslash}p{3cm}|}
\hline
\textbf{№} & \textbf{Вимога з безпеки інформації} & \textbf{Вимога ГПБ} & \multicolumn{2}{c|}{\textbf{ЦПБ}} \\
\cline{4-5}
& & & \textbf{Захід захисту} & \textbf{Налаштований зміст заходу захисту}\\
\hline
\endfirsthead

\hline
\textbf{№} & \textbf{Вимога з безпеки інформації} & \textbf{Вимога ГПБ} & \multicolumn{2}{c|}{\textbf{ЦПБ}} \\
\cline{4-5}
& & & \textbf{Захід захисту} & \textbf{Налаштований зміст заходу захисту}\\
\hline
\endhead

\hline
\endfoot

\hline
\endlastfoot

\multicolumn{5}{|>{\raggedright\arraybackslash}p{16cm}|}{\textbf{1. ІДЕНТИФІКАЦІЯ ТА АВТЕНТИФІКАЦІЯ (IA)}} \\* \noalign{\hrule height 0.4pt}
1 & ІДЕНТИФІКАЦІЯ ТА АВТЕНТИФІКАЦІЯ ПРИСТРОЇВ (IA-3) &  & IA-3 & Апаратне збереження ключів (Hardware-backed Keystore / Secure Enclave) та біометрична автентифікація:\newline - Приватні ключі (наприклад, secp384r1 / P256) зберігаються виключно в ізольованому апаратному середовищі: Secure Enclave для iOS/macOS (через CryptoKit) та Hardware-backed Keystore для Android.\newline - Ключовий матеріал генерується з атрибутом неекспортованості та захищається біометричною автентифікацією (TouchID/FaceID/BiometricPrompt).\newline - Dart-шар отримує лише абстрактний дескриптор (alias) та виконує криптографічні операції (наприклад, підпис) через захищений MethodChannel (x509\_\allowbreak{}multicast/keys), не отримуючи прямого доступу до приватного ключа.\vspace{1mm}\newline\noindent\rule{\linewidth}{0.4pt}\vspace{1mm}
Параметр: Параметри автентифікації пристроїв\newline Тип: string\newline Значення: \textbf{flutter\_\allowbreak{}v1} \\ \hline
\multicolumn{5}{|>{\raggedright\arraybackslash}p{16cm}|}{\textbf{2. ЗАХИСТ СИСТЕМ ТА КОМУНІКАЦІЙ (SC)}} \\* \noalign{\hrule height 0.4pt}
2 & КОНФІДЕНЦІЙНІСТЬ ТА ЦІЛІСНІСТЬ ПЕРЕДАЧІ (SC-8) & Забезпечити [Вибір (один або кілька): конфіденційність; цілісність] інформації, що передається. & SC-8 & Мінімізація витоків пам'яті та транзит даних через Platform Channels:\newline - Забезпечується абсолютно децентралізована P2P UDP Multicast комунікація без посередників/брокерів (Layer 1 Skynet).\newline - Платформні канали функціонують лише як транзит для \textquotedbl{}чорних скриньок\textquotedbl{} (зашифрованих CMS пакетів), унеможливлюючи появу відкритих текстів чи криптографічних ключів у платформо-залежному Swift/Kotlin коді.\newline - UI/Dart-шар (lib/main.dart) використовує метод \_\allowbreak{}sendEncryptedMessage для передачі зашифрованих даних як Uint8List (наприклад, після локального шифрування у CMS через FFI з Rust або C). Для тестування потоку додано кнопку \textquotedbl{}Send CMS Payload\textquotedbl{}.\newline - Нативний iOS шар (ios/Runner/AppDelegate.swift) та macOS шар (macos/Runner/MainFlutterWindow.swift) мають обробник sendEncryptedMessage, який очікує сирі байти (FlutterStandardTypedData) та передає їх безпосередньо у мережевий UDP-шар (MulticastService.shared.send(data:)), уникаючи розшифрування чи зберігання ключів у нативному коді.\vspace{1mm}\newline\noindent\rule{\linewidth}{0.4pt}\vspace{1mm}
Параметр: Алгоритми E2EE для комунікацій\newline Тип: string\newline Значення: \textbf{flutter\_\allowbreak{}v1} \\ \hline
3 & КРИПТОГРАФІЧНИЙ ЗАХИСТ (SC-13) & a. Визначити [Призначення: використання криптографічних засобів, визначених організацією];\newline b. Впровадити [Завдання: визначені організацією види криптографії для кожного визначеного криптографічного використання]. & SC-13 & Локальне криптографічне шифрування та формування CMS-конвертів:\newline - Уся обробка відкритих ключів або секретних даних (наприклад, шифрування ChaCha20-Poly1305 та формування CMS-конвертів) залишається інкапсульованою на обраному боці (наприклад, в Dart FFI модулі на базі Rust або C).\newline - Це усуває вразливість до перехоплення через MethodChannel, оскільки в платформні канали передаються виключно зашифровані пакети.\vspace{1mm}\newline\noindent\rule{\linewidth}{0.4pt}\vspace{1mm}
Параметр: Криптографічний провайдер\newline Тип: string\newline Значення: \textbf{flutter\_\allowbreak{}v1} \\ \hline
4 & АВТЕНТИФІКАЦІЯ СЕСІЇ (SC-23) & Забезпечити автентифікацію сеансів зв'язку. & SC-23 & Автентифікація сесій відбувається через перевірку електронних підписів в повідомленнях. Інфраструктура PKI (X.509) гарантує неспростовність сесій. Зв'язок Dart <-> Native захищений від перехоплення в межах процесу.\vspace{1mm}\newline\noindent\rule{\linewidth}{0.4pt}\vspace{1mm}
Параметр: Механізм автентифікації сесії\newline Тип: string\newline Значення: \textbf{flutter\_\allowbreak{}v1} \\ \hline
\multicolumn{5}{|>{\raggedright\arraybackslash}p{16cm}|}{\textbf{3. ПРИДБАННЯ СИСТЕМ ТА ПОСЛУГ (SA)}} \\* \noalign{\hrule height 0.4pt}
5 & ТЕСТУВАННЯ ВРАЗЛИВОСТЕЙ (SA-11) & Вимагати від розробника системи, системного компонента або системної служби на всіх етапах проєктування та життєвого циклу розробки системи:\newline a. створити та впровадити план з оцінювання безпеки та приватності;\newline b. виконати [Вибір (один або кілька): одиниця; інтеграція; система; регресія] тестування/оцінювання [Призначення: з визначеною організацією частотою] з [Призначення: визначена організацією глибиною та охопленням];\newline c. надати докази (свідчення) виконання плану оцінювання та результати тестування й оцінювань;\newline d. впровадити перевірку процесу виправлення недоліків;\newline e. виправити дефекти, виявлені під час тестування та оцінювання. & SA-11 & RASP (Runtime Application Self-Protection), виявлення скомпрометованих середовищ та обфускація:\newline - Впроваджено модуль RaspManager (на базі freerasp), який автоматично моніторить та реагує на запуск у Jailbroken/Rooted середовищах, використання емуляторів, підключення фреймворків для хукінгу (Frida, Xposed) та спроби перепакування чи підміни рушія (reFlutter). В разі виявлення загрози процес екстрено завершується.\newline - На рівні конфігурації підтримується Apple App Attest для iOS та Google Play Integrity для Android.\newline - Розроблено скрипт релізної збірки scripts/build\_\allowbreak{}release.sh. Збірка виконується з обов'язковими прапорцями --obfuscate та --split-debug-info, що ускладнює реверс-інжиніринг бізнес-логіки та приховує внутрішні класи і функції застосунку.\vspace{1mm}\newline\noindent\rule{\linewidth}{0.4pt}\vspace{1mm}
Параметр: Механізми тестування та самозахисту\newline Тип: string\newline Значення: \textbf{flutter\_\allowbreak{}v1} \\ \hline


\end{longtable}

\end{document}
